% Revision 1 abstract (P-22): DRAFT prepared by the assistant for the author (D-D).
% Plain text only (the Word builder accepts $d_z{=}x$ and \texttt{} and nothing else).
% [R: ...] marks values filled from cohort B2. Journal range: 150-200 words.
Choosing a post-hoc explainer for a tabular model trades faithfulness, stability and
computational cost, yet comparisons under matched conditions are scarce. This paper reports
a paired benchmark of SHAP, LIME, Anchors and DiCE in which every explainer sees the same
model, test instances and seed on one recorded host: five model families on the Adult
dataset (confirmatory), six further datasets and a configuration sweep (exploratory). Cost
is measured in time and in model evaluations, which do not depend on the hardware. On Adult,
SHAP's fidelity exceeds that of each alternative in all 25 model-seed cells ($d_z{=}3.3$ or
more), but its stability advantage over LIME and Anchors is not significant; DiCE ranks last
on every quality endpoint. On the six further datasets SHAP keeps the highest fidelity,
while LIME is the more stable explainer, as it is on Adult with category-preserving
noise or a wider kernel; SHAP's fidelity lead holds over every configuration tested. KernelSHAP's median wall time is 3 to 25 times LIME's,
whereas TreeSHAP's cost follows tree depth: faster than LIME on XGBoost, and on the deeper
random forest faster or slower by dataset. The choice of explainer therefore depends on the model, the
configuration and the latency budget.
